\section{Introduction}
\label{sec:intro}

The neural tangent kernel (NTK) is defined as $\Theta(x,x')=J(x)J(x')^{\top}$, where $J(x)$ is the Jacobian of the neural network evaluated at $x$ with respect to the network parameters. Given a training dataset and loss function, the NTK completely specifies the evolution of the neural network to first order in the learning rate. This makes the NTK a powerful tool to approximate neural network training dynamics and many works have measured $\Theta$ during or after training for various applications~\cite{hayase2022, mok2022, engel2023a, zhou2024, wilson2025}. At the same time, the NTK can also be studied theoretically. However, the NTK is a random variable via its dependence on the initialization and it evolves non-linearly during training as the parameters are updated, making direct theoretical analyses of the NTK prohibitively difficult.

Nevertheless, it is well-known that neural networks simplify considerably in the infinite-width limit, in which the number of channels of all hidden layers are taken to infinity~\cite{neal1996, lee2018, matthews2018}. In particular, the NTK collapses onto its mean and does not evolve during training, it \emph{freezes}~\cite{jacot2018}. This deterministic NTK can be computed in closed form for deep, non-linear networks by applying layer-wise recursive relations which have been derived for all common neural network layers~\cite{novak2020a}. Therefore, both gradient descent~\cite{jacot2018} and Bayesian inference~\cite{he2020} can be solved in closed form at infinite width.

% since the NTK captures the non-linear evolution of neural networks during training, it also shares many of its complexities. Namely


%Also the NTK simplifies considerably in this limit~\cite{jacot2018}. At infinite width, the distribution of the NTK collapses onto its mean, rendering the NTK deterministic. Furthermore, (when taking the right parametrization for the preactivations) the NTK does not evolve in training time, it \emph{freezes}. Finally, this constant, deterministic NTK can be computed in closed form for deep, non-linear networks by applying layer-wise recursive relations which have been derived for all common neural network layers~\cite{novak2020a}. These simplifications are so powerful that the training for both gradient descent~\cite{jacot2018} and Bayesian inference~\cite{he2020} can be solved in closed form at infinite width, leading to numerous insights about the training dynamics~\cite{xiao2020, geiger2020a, franceschi2022, wang2022b, deoliveira2025}. Importantly, these results pertain to non-linear, deep neural networks.

Although the study of the infinite-width limit has been very successful, the simplifications of the dynamics which enabled this progress also constitute an important limitation of this approach. In particular, at infinite width, the network is effectively linearized in its parameters~\cite{lee2019a} and becomes a Gaussian process. Furthermore, only the last layer evolves during training, a property termed \emph{no feature learning}. Correspondingly, some studies found considerable deviations between the properties observed for finite-width networks and those predicted by NTKs at infinite width~\cite{NEURIPS2020_ad086f59, wenger2024, liu2025}.

A solution for this problem is to consider how deviations from the strict infinite-width limit affect the NTK statistics. In particular, one can expand the training dynamics in a Taylor series in $1/n$, where $n$ is the width of the hidden layers~\cite{yaida2020, roberts2022}. In this expansion, the leading term, $1/n\rightarrow0$, corresponds to the infinite width limit and the first correction introduces finite-width effects. Including these finite-width effects in the analysis greatly extends the power of the theoretical framework which now includes non-trivial NTK statistics, the evolution of the NTK and feature learning.

Since the infinite-width statistics are Gaussian, this constitutes an expansion of the parameter statistics around a Gaussian distribution. Such an expansion is very common in theoretical physics, in particular in quantum field theory (QFT), which describes the interactions of elementary particles in a statistical manner. In perturbative QFT, the \emph{action} (log-probability) of the theory is expanded around a leading Gaussian contribution which corresponds to non-interacting particles, while non-Gaussian corrections introduce interactions. Therefore, by employing methods from QFT, one can compute corrections to the infinite-width limit of neural network statistics. However, the algebraic manipulations involved in deriving these results are lengthy and usually phrased in a language unfamiliar to the machine-learning community. In contrast, when performing computations in particle physics, physicists mostly use \emph{Feynman diagrams} as an intuitive shorthand to greatly simplify the computational effort.

In this article, we introduce Feynman-like diagrams for the computation of finite-width corrections to NTKs of MLPs. We provide Feynman rules for straightforwardly computing all relevant neural network statistics involving preactivations, the NTK and related quantities. In this framework, computing expectation values reduces to drawing all diagrams compatible with the desired expectation value and translating them into algebraic expressions using the Feynman rules we provide.

%With the framework we provide, it is possible to straightforwardly compute all relevant neural network statistics involving preactivations, the NTK and related quantities.
In order to demonstrate the power of our framework, we investigate the stability of gradients as the depth of the network increases. At infinite width, a critical initialization variance can be derived, for which the preactivations and NTK do not grow exponentially in depth. It could be shown that the same critical initialization stabilizes the forward pass also at finite width using Feynman diagrams for preactivations~\cite{banta2024}. Using our NTK-Feynman diagrams, we can extend this result to the gradients at finite width. We also show that for scale-invariant activation functions such as ReLU, $\Theta(x,x)$ does not include finite-width corrections and the infinite-width statistics are exact.

We confirm our theoretical results with numerical experiments. By computing ensemble-averages of preactivations and Jacobians, we demonstrate the stability of the critical initialization for statistics involving the NTK. In the same way, we confirm our stability results for scale-invariant nonlinearities.

Our main contributions are as follows:
\begin{itemize}
\item We provide a set of rules for computing finite-width corrections to NTK statistics using Feynman diagrams, greatly simplifying the algebraic manipulations required. We verify that our rules result in the correct layer-wise recursion relations for the NTK and preactivation statistics at order $1/\text{width}$ and find no contradictions for selected checks at higher order.
\item We demonstrate the power of our formalism by deriving the recursion relation of the leading order correction to the mean of the NTK using our Feynman diagrams. We furthermore extend previous results on preactivation stability to arbitrary statistics of preactivations and the NTK. Finally, by showing that all $1/\text{width}$ corrections of $\Theta(x,x)$ vanish for scale-invariant nonlinearities, we prove that its infinite-width expression is exact in this case.
\item We validate our results experimentally by sampling ensembles of MLPs. In particular, we show that the initialization variance which stabilizes the infinite-width limit also stabilizes all statistics involving the NTK to all orders in $1/\text{width}$ by sampling 7 tensors below, at and above criticality and measuring their dependence on depth. We furthermore verify by sampling that $\Theta(x,x)$ receives no finite-width corrections for the scale-invariant activation functions ReLU and LeakyReLU but corrections are present for GeLU (not scale-invariant).
\end{itemize}

%%% Local Variables:
%%% mode: LaTeX
%%% TeX-master: "ntk_feynman"
%%% End:
